%% -*- coding: utf-8 -*-
\documentclass[12pt,pagesize,paper=192mm:108mm,landscape]{scrbook} 
%1920x1080 1280x720
\areaset[current]{192mm}{108mm}
\usepackage{calc}
\usepackage[T2A]{fontenc}
\usepackage[utf8]{inputenc}
\usepackage[english,russian]{babel}
\usepackage{microtype}
\usepackage{misccorr}
\usepackage{cmap}
%\usepackage[unicode=true]{hyperref}
\usepackage{graphicx}
\usepackage{amssymb}
\usepackage{amsmath}
%\usepackage{srcltx}
\usepackage{textcomp}
\usepackage{xspace}
%научные символы и смайлики \smiley \frownie
\usepackage{wasysym}
\usepackage{ccicons}
%перенос формул в тексте
\newcommand*{\hm}[1]{#1\nobreak\discretionary{}%
  {\hbox{$\mathsurround=0pt #1$}}{}}

\begin{document}
\begin{titlepage}
  \vspace*{-0.5em}
  \begin{center}    
    % \hspace*{3em}
    % \begin{minipage}[t]{3em}
    %   \includegraphics[width=\textwidth]{../BINP-logo}
    % \end{minipage}\hfill
    % \begin{minipage}{0.23\linewidth}
    % \includegraphics[width=\textwidth]{../NSU-logo}
    % \end{minipage}
    % \hfill
    % \hspace*{6em}

    % Кафедра теоретической физики физического факультета НГУ
    % \medskip

    % \Large
    % Профессор Грабовский А.\,В.
    % \smallskip

    % \huge
    % \textbf{Общая теория относительности}
    % \smallskip

    \vfill
    \Large
    Задание № 10
 \bigskip

    \normalsize
    \begin{minipage}{0.75\linewidth}
      \begin{enumerate}
      \item Оценить полный поток энергии звездного света, сравнить с
        потоком от Солнца. Оценить, во сколько раз вселенная должна
        быть старше, чтобы небо ночью светилось как днем, считая, что
        звезды светили все время жизни вселенной с той же плотностью
        светимости, что и сейчас.

	\item Найти скаляр кривизны для трехмерной метрики \[dl^2 = a^2\left(dr^2/\left(1 - kr^2\right) + r^2d\Omega^2\right),\] \(k = +1,\, -1,\, 0\).
      \end{enumerate}
    \end{minipage}
    \vfill

     \normalsize \ccbysa\hspace{0.5em}  Новосибирск 2022
  \end{center}
\end{titlepage}
\end{document}

